\chapter{Conclusioni}
\label{ch:conclusioni}

Il lavoro ha confrontato, sullo stesso dominio HR/retributivo e sulla stessa infrastruttura, un relazionale
distribuito (PostgreSQL+Citus) e un documentale (MongoDB), entrambi CP ma con modelli di dati opposti. Ne
emerge che la domanda ``quale sistema è migliore'' non ha risposta assoluta: dipende dall'operazione e dalla
scala, coerentemente con l'idea che non esista una taglia unica~\cite{stonebraker2005onesize}.

\section{Sintesi dei risultati}
\label{sec:sintesi}

\begin{itemize}
  \item \textbf{Scalabilità per nodi (Citus)}: le analitiche cross-shard scalano quasi linearmente
        (${\sim}3.8\text{--}4.6\times$ con 4 nodi); le query mirate a un tenant restano piatte perché
        colpiscono un solo shard, ma sono già nell'ordine del millisecondo. Il miglioramento è più marcato
        proprio per le query che aggregano e uniscono (join) dati di più tenant, come mostrano
        \Cref{tab:scaling-nodi,tab:scaling-nodi-full}: la L05, ad esempio, arriva a $4.4\times$ passando da 1
        a 4 nodi.
  \item \textbf{Scalabilità per dati (Citus)}: a nodi fissi le cross-shard degradano super-linearmente col
        volume, mentre le single-shard restano insensibili grazie all'indicizzazione e alla co-locazione per
        tenant.
  \item \textbf{Scritture (Citus)}: la co-locazione fa sì che la scrittura di un tenant colpisca un solo nodo;
        solo la modifica di un dato di riferimento tocca tutti i nodi. La transazione multi-tabella è
        l'operazione più costosa.
  \item \textbf{Guasti}: comportamento \textbf{CP} confermato; con un nodo giù la query cross-shard
        fallisce per l'intera finestra (nessun risultato parziale), con recupero automatico e rapido al
        rientro. La consistenza è scelta al posto della disponibilità.
  \item \textbf{MongoDB}: più veloce sulle analitiche che l'embedding serve direttamente (percentili, spesa
        per categoria), ma \textbf{nettamente più lento} sulle query che attraversano più entità (assenteismo,
        ferie residue, che toccano assenze e ratei): lì Citus, con i join indicizzati, è risultato più rapido
        di 6--17 volte a pari scala. Inoltre il balancer, volutamente conservativo, \textbf{non distribuisce a
        scala piccola}, quindi a quella scala aggiungere shard non porta beneficio; la distribuzione riparte
        solo a volume grande, in contrasto con il rebalance esplicito e deterministico di Citus.
  \item \textbf{Modellazione}: il cedolino completo premia il documentale (un documento vs sette tabelle),
        mentre il dato di riferimento e le aggregazioni multi-entità premiano il relazionale (una riga vs
        aggiornamento di massa; join indicizzati vs scan di array annidati).
\end{itemize}

\section{Quando conviene ciascun sistema}
\label{sec:quando}

Sulla base delle misure: \textbf{Citus} è preferibile quando servono analitiche cross-tenant che devono
\emph{scalare} orizzontalmente in modo prevedibile e, soprattutto, per le \textbf{query complesse che
attraversano più entità} (aggregazioni e join profondi su assenze, ratei, contratti): proprio lì, dove
MongoDB deve scorrere gli array annidati sull'intera collezione, Citus è risultato molte volte più veloce a
pari scala. Si aggiungono l'integrità referenziale forte, le transazioni distribuite e il controllo esplicito
della distribuzione tipico di un DBA. \textbf{MongoDB} è preferibile quando il
carico è dominato da operazioni sull'\emph{aggregato} (lettura/scrittura del documento completo), quando conta
la bassa latenza sulle singole richieste e la flessibilità di schema, accettando in cambio un controllo meno
fine sulla distribuzione e la denormalizzazione dei dati condivisi.

\section{Limiti e sviluppi futuri}
\label{sec:limiti}

I risultati vanno letti alla luce di alcuni limiti dichiarati. La scala è dell'ordine dei GB (vincolo del
disco delle VM), non delle centinaia di GB. Il working set è sempre rimasto in RAM (cache di Postgres e di
WiredTiger): le letture da disco sono state nulle, quindi il regime I/O-bound non è stato raggiunto. Gli
esperimenti si sono fermati quando la tendenza era già chiara e i dati risultavano distribuiti su tutti i
nodi, senza spingersi fino a saturare il disco. Entrambi i sistemi sono inoltre a copia singola (nessuna
replica), scelta voluta per un test di guasto CP corretto ma che esclude l'alta disponibilità.

Come sviluppi: spingere i volumi fino alla saturazione effettiva del disco per osservare il regime I/O-bound;
introdurre la replica per confrontare disponibilità e comportamento in failover; e misurare esplicitamente
uno scenario di evoluzione di schema (turni/ferie) su entrambi i sistemi, oggi trattato solo a livello teorico.
